% Alluvial / Sankey flow of a single 511 keV annihilation photon's fate.
% \input by chapters/01-introduction.tex.
%
% The figure is QUALITATIVE and deliberately quotes no per-scanner numbers.
% Band widths are schematic: the trunk halves at every interface, so the scale
% is 0.06 cm per unit of a 100-unit trunk and the whole drawing is 6.40 cm
% tall. Nothing is measured off them. What the figure claims is in the ledger:
% which properties move each surviving fraction, which way, and roughly what
% range the fraction spans across the designs in use. The ranges are
% order-of-magnitude and still want a source.
%
% Two orderings the ledger must not contradict: body scatter and solid angle
% dominate the last two stages, and of those two only solid angle is a design
% choice.
%
% NOTE ON SCATTER: the first loss dead-ends here, but those photons are NOT
% absorbed. At 511 keV in tissue photoelectric absorption is negligible, so
% essentially all of them Compton scatter and most escape the body; a good
% share reach the detector and pass the energy window (a 435-650 keV window
% accepts scatter up to ~35 deg, since E' = 511/(2-cos(theta))), and are
% recorded as scattered coincidences. What the scattering destroys is the LOR,
% not the photon, and the trunk here tracks photons that keep a valid one.
%
% NOTE ON PAIRS: this tracks ONE photon. A coincidence needs both, so the
% per-photon terms square -- but the solid-angle term does NOT, since the
% photons are back-to-back and a ring that clears one largely clears the other.
% Axial length therefore pays once where crystal thickness pays twice.
%
% LAYOUT: five columns at x = 0, 2.2, 4.4, 6.6, 8.8, node bars 0.16 wide, a
% 0.40 cm gap between nodes sharing a column. Each loss terminates at the stage
% that claims it; only the trunk carries on, which is what stops the fat early
% losses from filling the canvas. Ribbons are cubic beziers with horizontal
% tangents (a pure vertical shear), so they cannot fold however thick they get,
% unlike the swept arcs the CTAN `sankey` package draws.

\providecolor{flowkeep}{HTML}{2A6F97}
\providecolor{flownode}{HTML}{1F4E6B}
\providecolor{flowloss1}{HTML}{CDD6DB}
\providecolor{flowloss2}{HTML}{AAB8C0}
\providecolor{flowloss3}{HTML}{8B9BA4}
\providecolor{flowloss4}{HTML}{64757E}
\providecolor{flowtext}{HTML}{3D4A52}

\begin{tikzpicture}[x=1cm, y=1cm, font=\footnotesize,
	ribbon/.style={fill opacity=0.75},
	node/.style={}]
\fill[flowkeep, ribbon]
	(0.080,0.000) .. controls (1.100,0.000) and (1.100,0.000) .. (2.120,0.000) -- (2.120,-3.000) .. controls (1.100,-3.000) and (1.100,-3.000) .. (0.080,-3.000) -- cycle;
\fill[flowloss1, ribbon]
	(0.080,-3.000) .. controls (1.100,-3.000) and (1.100,-3.400) .. (2.120,-3.400) -- (2.120,-6.400) .. controls (1.100,-6.400) and (1.100,-6.000) .. (0.080,-6.000) -- cycle;
\fill[flowkeep, ribbon]
	(2.280,0.000) .. controls (3.300,0.000) and (3.300,0.000) .. (4.320,0.000) -- (4.320,-1.500) .. controls (3.300,-1.500) and (3.300,-1.500) .. (2.280,-1.500) -- cycle;
\fill[flowloss2, ribbon]
	(2.280,-1.500) .. controls (3.300,-1.500) and (3.300,-1.900) .. (4.320,-1.900) -- (4.320,-3.400) .. controls (3.300,-3.400) and (3.300,-3.000) .. (2.280,-3.000) -- cycle;
\fill[flowkeep, ribbon]
	(4.480,0.000) .. controls (5.500,0.000) and (5.500,0.000) .. (6.520,0.000) -- (6.520,-0.750) .. controls (5.500,-0.750) and (5.500,-0.750) .. (4.480,-0.750) -- cycle;
\fill[flowloss3, ribbon]
	(4.480,-0.750) .. controls (5.500,-0.750) and (5.500,-1.150) .. (6.520,-1.150) -- (6.520,-1.900) .. controls (5.500,-1.900) and (5.500,-1.500) .. (4.480,-1.500) -- cycle;
\fill[flowkeep, ribbon]
	(6.680,0.000) .. controls (7.700,0.000) and (7.700,0.000) .. (8.720,0.000) -- (8.720,-0.375) .. controls (7.700,-0.375) and (7.700,-0.375) .. (6.680,-0.375) -- cycle;
\fill[flowloss4, ribbon]
	(6.680,-0.375) .. controls (7.700,-0.375) and (7.700,-0.775) .. (8.720,-0.775) -- (8.720,-1.150) .. controls (7.700,-1.150) and (7.700,-0.750) .. (6.680,-0.750) -- cycle;
\fill[flownode, node]
	(-0.080,0.000) rectangle (0.080,-6.000);
\fill[flowkeep, node]
	(2.120,0.000) rectangle (2.280,-3.000);
\fill[flowloss1, node]
	(2.120,-3.400) rectangle (2.280,-6.400);
\fill[flowkeep, node]
	(4.320,0.000) rectangle (4.480,-1.500);
\fill[flowloss2, node]
	(4.320,-1.900) rectangle (4.480,-3.400);
\fill[flowkeep, node]
	(6.520,0.000) rectangle (6.680,-0.750);
\fill[flowloss3, node]
	(6.520,-1.150) rectangle (6.680,-1.900);
\fill[flowkeep, node]
	(8.720,0.000) rectangle (8.880,-0.375);
\fill[flowloss4, node]
	(8.720,-0.775) rectangle (8.880,-1.150);

%% Stage labels. The cue under each is a pointer into the ledger, not a list --
%% neighbours are 2.2 cm apart and collide once two cues sum past 125 pt, which
%% produces no warning, so keep them near 45 pt each. Every \\ must sit at the
%% top level of the node text; one nested inside a {\scriptsize ...} group
%% breaks TikZ's line breaker, hence the bare font switch.
\node[anchor=south, align=center, text=flowtext] at (1.10,0.30)
	{Body\\transport\\[2pt]\scriptsize\itshape patient size\\and density};
\node[anchor=south, align=center, text=flowtext] at (3.30,0.30)
	{Solid\\angle\\[2pt]\scriptsize\itshape axial FOV and\\ring diameter};
\node[anchor=south, align=center, text=flowtext] at (5.50,0.30)
	{Crystal\\interaction\\[2pt]\scriptsize\itshape density, $Z_{\mathrm{eff}}$,\\thickness};
\node[anchor=south, align=center, text=flowtext] at (7.70,0.30)
	{Energy\\window\\[2pt]\scriptsize\itshape light yield and\\resolution};

\node[anchor=east, align=right, text=flowtext] at (-0.18,-3.000)
	{\keV{511}\\photons\\emitted};

%% Terminal losses, named beside their own node -- the ribbons are too narrow
%% at the waist to hold text, and every dependence lives in the ledger anyway.
\node[anchor=west, align=left, text=flowtext] at (2.35,-4.900)
	{Compton scattered in the body,\\so the line of response is lost};
\node[anchor=west, align=left, text=flowtext] at (4.55,-2.650)
	{Emitted outside the\\detector's solid angle};

%% The right-hand stack is the tightest spacing in the figure: three two-line
%% labels (0.76 cm each) sharing the top 2.6 cm, set to leave 0.2 cm of white
%% between blocks. The crystal leader must also stay above the solid-angle
%% label, whose first line ends near x = 7.05 and whose top edge is at -2.27;
%% the leader is at -1.63 there.
\draw[flowtext, line width=0.3pt] (8.88,-0.960) -- (8.94,-1.150);
\node[anchor=west, align=left, text=flowtext] at (9.00,-1.150)
	{Rejected by the\\energy window};
\draw[flowtext, line width=0.3pt] (6.70,-1.525) -- (8.94,-2.200);
\node[anchor=west, align=left, text=flowtext] at (9.00,-2.200)
	{Passed through with no\\interaction in the crystal};
\node[anchor=west, align=left, text=flowkeep] at (9.00,-0.188)
	{\textbf{Recorded single}\\\textbf{with a valid LOR}};

%% The ledger: one row per interface, direction of every dependence plus the
%% range it spans. 2.35+0.3+1.55+0.3+5.35+0.3+5.35 = 15.5 cm, centred on the
%% drawing's midpoint. The closing note is a separate node because a p-column
%% cell that opens or closes with its own \\ leaves the alignment mid-group and
%% TeX halts with "Missing \endgroup inserted"; keep prose out of the tabular.
\node[anchor=north, text=flowtext, font=\scriptsize] (ledger) at (5.25,-6.80) {%
	\begin{tabular}{@{}p{2.35cm} @{\hspace{0.3cm}} p{1.55cm} @{\hspace{0.3cm}} p{5.35cm} @{\hspace{0.3cm}} p{5.35cm}@{}}
		\toprule
		\textbf{Interface} & \textbf{Range} & \textbf{Larger surviving fraction} & \textbf{Smaller surviving fraction} \\
		\midrule
		\textcolor{flowloss1}{\rule[-0.02em]{0.55em}{0.55em}}~Body transport
			& $\approx$\,15--60\,\%
			& Smaller or less dense patient; source nearer the surface.
			& Large or dense patient; deep-lying source. Set by the patient, not the scanner. \\
		\addlinespace
		\textcolor{flowloss2}{\rule[-0.02em]{0.55em}{0.55em}}~Solid angle
			& $\approx$\,20--80\,\%
			& Longer axial field of view; smaller ring diameter; source near the axial centre.
			& Short axial extent; wide bore; source near the end of the field of view. \\
		\addlinespace
		\textcolor{flowloss3}{\rule[-0.02em]{0.55em}{0.55em}}~Crystal interaction
			& $\approx$\,60--90\,\%
			& Denser, higher-$Z_{\mathrm{eff}}$, thicker scintillator; high packing fraction, so little of the face is reflector or gap.
			& Thin or low-density crystals; thick reflector between elements. Thickness is bought against depth-of-interaction blur. \\
		\addlinespace
		\textcolor{flowloss4}{\rule[-0.02em]{0.55em}{0.55em}}~Energy window
			& $\approx$\,40--70\,\%
			& High and uniform light yield, hence better energy resolution; high $Z_{\mathrm{eff}}$; a wider window.
			& Poor or non-uniform light collection; low $Z_{\mathrm{eff}}$; a narrow window. \\
		\bottomrule
	\end{tabular}};
\node[anchor=north west, text=flowtext, font=\scriptsize\itshape,
	text width=15.5cm, align=left] at ([yshift=-3pt]ledger.south west)
	{Each range is the fraction surviving that interface alone, not of all emitted photons, and spans the designs and patients in ordinary use rather than the extremes. The last row is the only one where a larger fraction is not automatically better, since the window widens to admit scatter as well as photopeak events; note too that the photofraction and the fraction reaching the photopeak differ, because Compton interactions that deposit their remaining energy in the same crystal still land in the photopeak.\needcite{ranges for each of the four interfaces -- a review, or the author's own Monte Carlo model}};
\end{tikzpicture}
